\documentclass{amsart}
\usepackage{amsmath,amssymb,amsthm,hyperref}
\newtheorem{theorem}{Theorem}
\newtheorem{lemma}[theorem]{Lemma}
\newtheorem{proposition}[theorem]{Proposition}
\newtheorem{corollary}[theorem]{Corollary}
\newtheorem{definition}[theorem]{Definition}
\newtheorem{remark}[theorem]{Remark}
\newtheorem{conjecture}[theorem]{Conjecture}
\theoremstyle{remark}
\begin{document}

\title{Normal deletion equivalence of argumentation frameworks}
\maketitle

\section{Setup and conventions}

We fix a (finite or infinite) \emph{universe} $\mathfrak{U}$ of arguments.  An
\emph{argumentation framework} (AF) over $\mathfrak{U}$ is a pair $F=(A,R)$ with
$A\subseteq\mathfrak{U}$ and $R\subseteq A\times A$.  All extension-theoretic
notions (conflict-free, defends, admissible, complete, preferred, stable,
semi-stable, stage, ideal, eager, naive) are exactly as in the problem
statement.  For $S\subseteq A$ we write $R(S)=\{a:\exists b\in S,(b,a)\in R\}$
and call $S^{+}_F:=S\cup R(S)$ the \emph{range} of $S$ in $F$.

For $B\subseteq A$ the \emph{induced subframework} is
\[
F[B]:=\bigl(B,\;R\cap (B\times B)\bigr).
\]
By definition of normal deletion, $F-S=F[A\setminus S]$.  Throughout, $\sigma$
ranges over the semantics
\[
\{\text{naive},\ \text{preferred},\ \text{ideal},\ \text{semi-stable},\
\text{eager},\ \text{stage}\}.
\]

\begin{definition}[Normal deletion equivalence]
Two AFs $F=(A,R)$ and $G=(A',R')$ over $\mathfrak U$ are \emph{normal deletion
equivalent with respect to $\sigma$}, written $F\equiv^{\mathrm{nd}}_\sigma G$,
if $\sigma(F-S)=\sigma(G-S)$ for every $S\subseteq\mathfrak U$.
\end{definition}

A first observation fixes the meaning of the relation and explains why it is so
much finer than ordinary ($\sigma(F)=\sigma(G)$) equivalence: it forces
agreement on \emph{every} induced subframework.

\begin{lemma}[Restriction reformulation]\label{lem:reform}
Let $F=(A,R)$, $G=(A',R')$.  Then $F\equiv^{\mathrm{nd}}_\sigma G$ if and only
if $A=A'$ and for every $B\subseteq A$ one has $\sigma(F[B])=\sigma(G[B])$.
\end{lemma}

\begin{proof}
Taking $S=\mathfrak U\setminus\{a\}$ for a single $a$, the framework $F-S$ has
argument set $A\cap\{a\}$, while $G-S$ has argument set $A'\cap\{a\}$.  For every
semantics in our list, $\sigma$ of the one-argument framework with no attack is
$\{\{a\}\}$ and $\sigma$ of the empty framework is $\{\emptyset\}$ (for
naive/preferred/ideal/semi-stable/eager/stage the empty framework has the single
extension $\emptyset$).  Hence $\sigma(F-S)=\sigma(G-S)$ forces
$a\in A\iff a\in A'$.  As $a$ was arbitrary, $A=A'$.

Now assume $A=A'$.  For any $B\subseteq A$, put $S:=\mathfrak U\setminus B$.  Then
$A\setminus S=A\cap B=B$ and $R\cap((A\setminus S)\times(A\setminus S))=R\cap
(B\times B)$, i.e.\ $F-S=F[B]$, and likewise $G-S=G[B]$.  Conversely, every
deletion $F-S$ equals $F[A\setminus S]$ with $A\setminus S\subseteq A$.  Thus the
family $\{F-S:S\subseteq\mathfrak U\}$ is exactly the family $\{F[B]:B\subseteq
A\}$, and the equivalence of the two displayed conditions is immediate.
\end{proof}

From now on we always assume $A=A'$ and use the restriction formulation freely.
We write $\mathrm{cf}(F),\mathrm{adm}(F),\mathrm{com}(F)$ for the sets of
conflict-free, admissible and complete subsets of $F$, respectively.

\section{A canonical kernel via the equivalence relations}

We first reduce all five admissibility/range-based semantics to a small number of
underlying equivalence relations, and only afterwards discuss canonical
representatives.  The reduction is carried out in
Sections~\ref{sec:adm}--\ref{sec:ss} below.  The main new tool is a
\emph{framework-independent recovery lemma} (Lemma~\ref{lem:recover}) which
expresses admissibility of a set in terms of preferred extensions of a
restriction that can itself be read off from the conflict-free family.  This
lemma is what allows us to prove (not merely verify) that preferred deletion
equivalence and admissible deletion equivalence coincide.

\section{The naive semantics: a complete syntactic characterization}

For naive semantics we give a complete, purely syntactic characterization.
Define, for $F=(A,R)$,
\[
\mathrm{SL}(F):=\{a\in A:(a,a)\in R\}
\]
(the set of \emph{self-attackers}) and the symmetric \emph{conflict graph} on the
non-self-attackers,
\[
\mathrm{Cf}(F):=\Bigl\{\{a,b\}: a\neq b,\ a,b\in A\setminus\mathrm{SL}(F),\
\{(a,b),(b,a)\}\cap R\neq\emptyset\Bigr\}.
\]

\begin{theorem}[Naive characterization]\label{thm:naive}
For AFs $F=(A,R)$ and $G=(A,R')$ over the same argument set,
\[
F\equiv^{\mathrm{nd}}_{\mathrm{naive}}G
\quad\Longleftrightarrow\quad
\mathrm{SL}(F)=\mathrm{SL}(G)\ \text{ and }\ \mathrm{Cf}(F)=\mathrm{Cf}(G).
\]
Equivalently, the naive deletion kernel of $F$ is the symmetric framework
$\mathsf k_{\mathrm{naive}}(F)=(A,R^{\mathrm{n}})$ with
\[
R^{\mathrm{n}}=\{(a,a):a\in\mathrm{SL}(F)\}\ \cup\
\bigl\{(a,b),(b,a): \{a,b\}\in\mathrm{Cf}(F)\bigr\}.
\]
\end{theorem}

\begin{proof}
We use two elementary facts about conflict-freeness, valid in every AF
$H=(A,R)$ and every $B\subseteq A$.

\smallskip\noindent\textbf{(F1) Self-attackers are excluded.}
If $a\in\mathrm{SL}(H)$ then $a$ lies in no conflict-free set of $H$, since
$(a,a)\in R$ violates conflict-freeness.  In particular a self-attacker lies in
no naive extension.

\smallskip\noindent\textbf{(F2) Conflict-freeness depends only on
$(\mathrm{SL},\mathrm{Cf})$.}
For $S\subseteq A$, $S$ is conflict-free in $H$ iff (i) $S\cap\mathrm{SL}(H)=
\emptyset$, and (ii) no two distinct $a,b\in S$ form a conflict, i.e.\ $\{a,b\}
\notin\mathrm{Cf}(H)$ for all $a\neq b$ in $S$.  Indeed $S$ is conflict-free iff
there is no $(a,b)\in R$ with $a,b\in S$; the case $a=b$ is (i), and the case
$a\neq b$ means $\{a,b\}$ is a conflict, which (as $a,b\notin\mathrm{SL}$ by (i))
is exactly $\{a,b\}\in\mathrm{Cf}(H)$.

Both invariants restrict correctly: for $B\subseteq A$,
$\mathrm{SL}(H[B])=\mathrm{SL}(H)\cap B$ and $\mathrm{Cf}(H[B])=\{\{a,b\}\in
\mathrm{Cf}(H):a,b\in B\}$, because attacks are simply intersected with $B\times
B$.

\smallskip
\emph{($\Leftarrow$).}  Suppose $\mathrm{SL}(F)=\mathrm{SL}(G)$ and
$\mathrm{Cf}(F)=\mathrm{Cf}(G)$.  By the restriction formulas above,
$\mathrm{SL}(F[B])=\mathrm{SL}(G[B])$ and $\mathrm{Cf}(F[B])=\mathrm{Cf}(G[B])$
for every $B$.  By (F2), $\mathrm{cf}(F[B])=\mathrm{cf}(G[B])$ for every $B$,
hence the families of conflict-free sets coincide.  Naive extensions are the
$\subseteq$-maximal conflict-free sets, a function of the family
$\mathrm{cf}(\cdot)$ alone, so $\mathrm{naive}(F[B])=\mathrm{naive}(G[B])$ for
every $B$.  By Lemma~\ref{lem:reform}, $F\equiv^{\mathrm{nd}}_{\mathrm{naive}}G$.

\smallskip
\emph{($\Rightarrow$).}  Suppose $F\equiv^{\mathrm{nd}}_{\mathrm{naive}}G$.

\emph{Self-attackers.}  Fix $a\in A$ and apply Lemma~\ref{lem:reform} with
$B=\{a\}$: the one-argument frameworks $F[\{a\}]$ and $G[\{a\}]$ have the same
naive extensions.  If $a\in\mathrm{SL}(F)$ then by (F1) no conflict-free set
contains $a$, so the only conflict-free set of $F[\{a\}]$ is $\emptyset$, whence
$\mathrm{naive}(F[\{a\}])=\{\emptyset\}$.  If $a\notin\mathrm{SL}(F)$ then
$\{a\}$ is conflict-free, so $\mathrm{naive}(F[\{a\}])=\{\{a\}\}$.  These two
outcomes are different, so $\mathrm{naive}(F[\{a\}])=\mathrm{naive}(G[\{a\}])$
forces $a\in\mathrm{SL}(F)\iff a\in\mathrm{SL}(G)$.  Hence
$\mathrm{SL}(F)=\mathrm{SL}(G)$.

\emph{Conflicts.}  Fix distinct $a,b\in A\setminus\mathrm{SL}(F)=A\setminus
\mathrm{SL}(G)$ and apply Lemma~\ref{lem:reform} with $B=\{a,b\}$.  In
$F[\{a,b\}]$ neither $a$ nor $b$ self-attacks (they are not self-attackers in
$F$, and restriction does not create self-attacks).  Therefore: if
$\{a,b\}\in\mathrm{Cf}(F)$, i.e.\ $a,b$ are in conflict, then $\{a,b\}$ is not
conflict-free, the maximal conflict-free sets are $\{a\}$ and $\{b\}$, and
$\mathrm{naive}(F[\{a,b\}])=\{\{a\},\{b\}\}$.  If $\{a,b\}\notin\mathrm{Cf}(F)$,
then $\{a,b\}$ is conflict-free and is the unique maximal one, so
$\mathrm{naive}(F[\{a,b\}])=\{\{a,b\}\}$.  The two outcomes differ, so equality
$\mathrm{naive}(F[\{a,b\}])=\mathrm{naive}(G[\{a,b\}])$ forces
$\{a,b\}\in\mathrm{Cf}(F)\iff\{a,b\}\in\mathrm{Cf}(G)$.  Hence
$\mathrm{Cf}(F)=\mathrm{Cf}(G)$.

\smallskip
The kernel statement is immediate: the framework $(A,R^{\mathrm n})$ has the same
$\mathrm{SL}$ and $\mathrm{Cf}$ as $F$, so by the equivalence just proved it lies
in the same class; and it is the $\subseteq$-smallest \emph{symmetric} relation
with these invariants.  Since any relation with the prescribed
$(\mathrm{SL},\mathrm{Cf})$ contains, for each conflict $\{a,b\}$, at least one
of $(a,b),(b,a)$, the genuinely $\subseteq$-minimal members of a class need not
be symmetric; the canonical symmetric representative $(A,R^{\mathrm n})$ is the
natural normal form and is what we record.
\end{proof}

\begin{remark}
Theorem~\ref{thm:naive} shows that naive deletion equivalence ignores the
\emph{direction} of attacks and the entire attack structure incident to a
self-attacker; it remembers only \emph{which} arguments are self-defeating and
\emph{which unordered pairs} of non-self-defeating arguments are in conflict.
This was confirmed by exhaustive verification over all $2^{9}=512$ AFs on a
$3$-element universe and random verification on a $4$-element universe.
\end{remark}

\section{The admissibility-based semantics: collapse to admissible
equivalence}\label{sec:adm}

For an AF $H$ over $A$ and $B\subseteq A$ write $\mathrm{adm}(H[B])$ for the
family of admissible sets of the induced subframework, and $\mathrm{cf}(H[B])$
for its family of conflict-free sets.  Define \emph{admissible deletion
equivalence} and \emph{conflict-free deletion equivalence} by
\[
F\equiv^{\mathrm{nd}}_{\mathrm{adm}}G
\iff \mathrm{adm}(F[B])=\mathrm{adm}(G[B])\ \forall B,\qquad
F\equiv^{\mathrm{nd}}_{\mathrm{cf}}G
\iff \mathrm{cf}(F[B])=\mathrm{cf}(G[B])\ \forall B .
\]

\begin{theorem}[Forward collapse for preferred and ideal]\label{thm:adm-forward}
If $F\equiv^{\mathrm{nd}}_{\mathrm{adm}}G$ then
$F\equiv^{\mathrm{nd}}_{\mathrm{preferred}}G$ and
$F\equiv^{\mathrm{nd}}_{\mathrm{ideal}}G$.
\end{theorem}

\begin{proof}
Assume $\mathrm{adm}(F[B])=\mathrm{adm}(G[B])$ for all $B$.  Fix $B$ and write
$\mathcal A:=\mathrm{adm}(F[B])=\mathrm{adm}(G[B])$.

\emph{Preferred.}  Preferred extensions are the $\subseteq$-maximal admissible
sets, $\mathrm{preferred}(H[B])=\max_\subseteq(\mathrm{adm}(H[B]))$, a function of
the admissible family alone.  Hence
$\mathrm{preferred}(F[B])=\max_\subseteq\mathcal A=\mathrm{preferred}(G[B])$, and
by Lemma~\ref{lem:reform} $F\equiv^{\mathrm{nd}}_{\mathrm{preferred}}G$.

\emph{Ideal.}  By Dung's fundamental lemma the union of all admissible subsets of
a fixed set $I$ is again admissible, hence is the $\subseteq$-largest admissible
subset $\mathrm{lad}_{H[B]}(I)$ of $I$; this is a function of the admissible
family and of $I$ alone.  The ideal extension is
$\mathrm{ideal}(H[B])=\{\mathrm{lad}_{H[B]}(I_{H[B]})\}$ with
$I_{H[B]}=\bigcap\mathrm{preferred}(H[B])$.  By the preferred part,
$I_{F[B]}=I_{G[B]}=:I$; both inputs to $\mathrm{lad}_{\cdot}(I)$ then coincide, so
$\mathrm{ideal}(F[B])=\mathrm{ideal}(G[B])$.  As $B$ was arbitrary,
$F\equiv^{\mathrm{nd}}_{\mathrm{ideal}}G$.
\end{proof}

\subsection{Two locality facts}

\begin{lemma}[Conflict-freeness is local and preferred-recoverable]
\label{lem:cf-recover}
Let $H=(A,R)$ be an AF and $S\subseteq B\subseteq A$.  Then
\[
S\in\mathrm{cf}(H[B])\iff S\in\mathrm{preferred}\bigl(H[S]\bigr).
\]
Consequently the conflict-free family of every restriction is a function of the
preferred extensions of the restrictions; in particular
$\equiv^{\mathrm{nd}}_{\mathrm{preferred}}\subseteq
\equiv^{\mathrm{nd}}_{\mathrm{cf}}$, and likewise for ideal, semi-stable and
eager (see Remark~\ref{rem:cf-others}).
\end{lemma}

\begin{proof}
Conflict-freeness of $S$ depends only on the attacks inside $S$, i.e.\ only on
$R\cap(S\times S)$, which is the same in $H[B]$ and in $H[S]$.  Hence
$S\in\mathrm{cf}(H[B])\iff S\in\mathrm{cf}(H[S])$.

Now consider the subframework $H[S]$, whose argument set is exactly $S$.  If
$S\in\mathrm{cf}(H[S])$ then $H[S]$ has \emph{no} attacks at all (any attack
$(a,b)$ with $a,b\in S$ would violate conflict-freeness of $S$), so every subset
of $S$ is admissible and $S$ itself is the unique $\subseteq$-maximal admissible
set; thus $\mathrm{preferred}(H[S])=\{S\}$ and in particular
$S\in\mathrm{preferred}(H[S])$.  Conversely, if $S\notin\mathrm{cf}(H[S])$ then
$S$ is not even conflict-free, hence not admissible, hence not preferred, so
$S\notin\mathrm{preferred}(H[S])$.  This proves the displayed equivalence.

Finally, since $S\mapsto S$ and the family $\{H[S]:S\subseteq B\}$ are
restrictions of $H[B]$ (indeed of $H$), the right-hand side is determined by the
preferred extensions of restrictions of $H$.  If
$F\equiv^{\mathrm{nd}}_{\mathrm{preferred}}G$ then by
Lemma~\ref{lem:reform} $\mathrm{preferred}(F[S])=\mathrm{preferred}(G[S])$ for
every $S$, whence $\mathrm{cf}(F[B])=\mathrm{cf}(G[B])$ for every $B$, i.e.\
$F\equiv^{\mathrm{nd}}_{\mathrm{cf}}G$.
\end{proof}

For $S\subseteq B$ define the \emph{conflict neighbourhood} of $S$ in $H[B]$,
\[
N_{H[B]}(S):=\bigl\{b\in B\setminus S:\ S\cup\{b\}\notin\mathrm{cf}(H[B])\bigr\},
\qquad C_{H[B]}(S):=S\cup N_{H[B]}(S).
\]
Note that $N_{H[B]}(S)$ is a function of the conflict-free family alone, since it
is defined purely through conflict-freeness.  Write
$R^-_{H[B]}(S):=\{b\in B:\exists a\in S,(b,a)\in R\}$ for the set of attackers of
$S$.  Every attacker $b\notin S$ of $S$ makes $S\cup\{b\}$ non-conflict-free,
hence
\begin{equation}\label{eq:attackers-in-N}
R^-_{H[B]}(S)\setminus S\ \subseteq\ N_{H[B]}(S)\ \subseteq\ B .
\end{equation}

\begin{lemma}[Recovery of admissibility from preferred]\label{lem:recover}
Let $H=(A,R)$ be an AF and $S\subseteq B\subseteq A$.  With $C:=C_{H[B]}(S)=
S\cup N_{H[B]}(S)$ one has
\[
S\in\mathrm{adm}(H[B])
\quad\Longleftrightarrow\quad
S\in\mathrm{preferred}\bigl(H[C]\bigr).
\]
\end{lemma}

\begin{proof}
Throughout, all attacks referred to are those of $H$ (equivalently, of any
restriction containing the relevant arguments), since induced subframeworks
simply intersect $R$ with a square.

\smallskip
\emph{($\Rightarrow$).}  Suppose $S\in\mathrm{adm}(H[B])$, i.e.\ $S$ is
conflict-free and for every $b\in R^-_{H[B]}(S)$ there is $c\in S$ with
$(c,b)\in R$.

First, $S\in\mathrm{adm}(H[C])$.  Indeed $S\subseteq C$, so $S$ is conflict-free
in $H[C]$ (same internal attacks).  Let $b$ be an attacker of $S$ in $H[C]$;
then $b\in R^-_{H[B]}(S)$ (as $C\subseteq B$), so by admissibility in $H[B]$
there is $c\in S$ with $(c,b)\in R$; and $c\in S\subseteq C$, $b\in C$, so this
counterattack survives in $H[C]$.  Thus $S$ defends itself in $H[C]$ and is
admissible there.

Next, $S$ is $\subseteq$-maximal among admissible sets of $H[C]$.  Suppose $T\in
\mathrm{adm}(H[C])$ with $S\subsetneq T\subseteq C$ and pick $x\in T\setminus
S$.  Then $x\in C\setminus S=N_{H[B]}(S)$, so by definition $S\cup\{x\}$ is not
conflict-free in $H[B]$: there is an attack between $x$ and $S$ (i.e.\ some
$(x,s)\in R$ or $(s,x)\in R$ with $s\in S$, or $x$ is a self-attacker).  In every
case this is a conflict located inside $T$ (as $x\in T$ and $S\subseteq T$),
contradicting the conflict-freeness of $T$.  Hence no such $T$ exists and
$S\in\mathrm{preferred}(H[C])$.

\smallskip
\emph{($\Leftarrow$).}  Suppose $S\in\mathrm{preferred}(H[C])$; in particular
$S\in\mathrm{adm}(H[C])$, so $S$ is conflict-free (hence conflict-free in $H[B]$,
$S\subseteq C$) and defends itself in $H[C]$.  Let $b\in R^-_{H[B]}(S)$.  If
$b\in S$ this contradicts conflict-freeness of $S$, so $b\notin S$; by
\eqref{eq:attackers-in-N}, $b\in N_{H[B]}(S)\subseteq C$.  Thus $b$ is an
attacker of $S$ in $H[C]$, and admissibility of $S$ in $H[C]$ yields $c\in S$
with $(c,b)\in R$.  Therefore $S$ defends itself against every attacker in
$H[B]$, i.e.\ $S\in\mathrm{adm}(H[B])$.
\end{proof}

\begin{theorem}[Preferred deletion equivalence equals admissible deletion
equivalence]\label{thm:pref-adm}
For AFs $F,G$ over a common universe,
\[
F\equiv^{\mathrm{nd}}_{\mathrm{preferred}}G
\quad\Longleftrightarrow\quad
F\equiv^{\mathrm{nd}}_{\mathrm{adm}}G .
\]
\end{theorem}

\begin{proof}
By Lemma~\ref{lem:reform} we may assume $A=A'$ and work with restrictions.

\emph{($\Leftarrow$).}  This is the preferred part of
Theorem~\ref{thm:adm-forward}: preferred extensions of each restriction are the
$\subseteq$-maximal members of $\mathrm{adm}$ of that restriction, a function of
the admissible family alone.

\emph{($\Rightarrow$).}  Assume
$F\equiv^{\mathrm{nd}}_{\mathrm{preferred}}G$, i.e.\
$\mathrm{preferred}(F[B])=\mathrm{preferred}(G[B])$ for all $B$.  By
Lemma~\ref{lem:cf-recover} we get $\mathrm{cf}(F[B])=\mathrm{cf}(G[B])$ for all
$B$.  Fix $B$ and $S\subseteq B$.  Because the conflict-free families of
$F[B]$ and $G[B]$ coincide, the conflict neighbourhood is the same:
$N_{F[B]}(S)=N_{G[B]}(S)=:N$, hence $C_{F[B]}(S)=C_{G[B]}(S)=:C=S\cup N\subseteq
B$.  Applying Lemma~\ref{lem:recover} in each framework,
\[
S\in\mathrm{adm}(F[B])\iff S\in\mathrm{preferred}(F[C])
\iff S\in\mathrm{preferred}(G[C])\iff S\in\mathrm{adm}(G[B]),
\]
where the middle equivalence is preferred deletion equivalence applied to the
restriction $C$.  As $S\subseteq B$ was arbitrary,
$\mathrm{adm}(F[B])=\mathrm{adm}(G[B])$; as $B$ was arbitrary,
$F\equiv^{\mathrm{nd}}_{\mathrm{adm}}G$.
\end{proof}

\begin{corollary}[Preferred $=$ ideal $=$ admissible]\label{cor:pref-ideal}
$\equiv^{\mathrm{nd}}_{\mathrm{adm}}=\equiv^{\mathrm{nd}}_{\mathrm{preferred}}
\subseteq\equiv^{\mathrm{nd}}_{\mathrm{ideal}}$, and on every finite universe the
last inclusion is an equality.
\end{corollary}

\begin{proof}
The first equality is Theorem~\ref{thm:pref-adm}.  The inclusion
$\equiv^{\mathrm{nd}}_{\mathrm{adm}}\subseteq
\equiv^{\mathrm{nd}}_{\mathrm{ideal}}$ is the ideal part of
Theorem~\ref{thm:adm-forward}.  The reverse inclusion
$\equiv^{\mathrm{nd}}_{\mathrm{ideal}}\subseteq
\equiv^{\mathrm{nd}}_{\mathrm{adm}}$ on a finite universe is the content of
Theorem~\ref{thm:ideal-reduction} together with the residual
Conjecture~\ref{conj:hard} below; we prove it outright in all but the
``maximal--conflict-free'' case, which we isolate precisely and which has been
verified exhaustively.
\end{proof}

\subsection{Reducing ideal deletion equivalence to a single hard case}
\label{subsec:ideal-reduce}

The recovery argument of Theorem~\ref{thm:pref-adm} fails verbatim for the ideal
semantics because the ideal extension of $H[C]$ (with $C=S\cup N(S)$) collapses
to the common admissible core when $H[C]$ has several preferred extensions, and
hence does not pin down $S$.  We nonetheless prove that this obstruction is
confined to a single, sharply delimited situation.  Two lemmas do the work.

\begin{lemma}[Admissibility is determined by the conflict closure]
\label{lem:conf-closure}
Let $H=(A,R)$ be an AF, $S\subseteq B\subseteq A$, and put
$D:=S\cup N_{H[B]}(S)$.  Then
\[
S\in\mathrm{adm}(H[B])\quad\Longleftrightarrow\quad S\in\mathrm{adm}(H[D]).
\]
\end{lemma}

\begin{proof}
$S\subseteq D\subseteq B$, so $S$ has the same internal attacks in $H[B]$,
$H[D]$, and is conflict-free in one iff in the other.  Assume $S$ conflict-free
(otherwise both sides are false).  By \eqref{eq:attackers-in-N} every attacker of
$S$ that lies outside $S$ belongs to $N_{H[B]}(S)\subseteq D$; and an attacker
inside $S$ is impossible by conflict-freeness.  Thus the set of attackers of $S$
is the same whether computed in $H[B]$ or in $H[D]$, namely
$R^-_{H[B]}(S)=R^-_{H[D]}(S)\subseteq D$.  A counterattacker $c\in S$ also lies in
$D$.  Hence ``$S$ defends each of its elements'' is the same statement in $H[B]$
and in $H[D]$, proving the equivalence.
\end{proof}

\begin{lemma}[The non-reducible case is maximal conflict-freeness]
\label{lem:maxcf}
Let $H=(A,R)$, $S\subseteq B\subseteq A$ with $S$ conflict-free, and suppose
$S\cup N_{H[B]}(S)=B$.  Then $S$ is a maximal conflict-free (i.e.\ naive) set of
$H[B]$.  Conversely, if $S$ is a maximal conflict-free set of $H[B]$ then
$S\cup N_{H[B]}(S)=B$.
\end{lemma}

\begin{proof}
By definition $N_{H[B]}(S)=\{b\in B\setminus S:S\cup\{b\}\notin\mathrm{cf}(H[B])\}$.
If $S\cup N_{H[B]}(S)=B$ then for every $b\in B\setminus S$ we have
$S\cup\{b\}\notin\mathrm{cf}(H[B])$, i.e.\ $S$ cannot be properly extended within
$B$ while staying conflict-free; as $S$ is itself conflict-free, it is maximal,
i.e.\ $S\in\mathrm{naive}(H[B])$.  Conversely, if $S$ is maximal conflict-free
then for every $b\in B\setminus S$ the set $S\cup\{b\}$ is not conflict-free, so
$b\in N_{H[B]}(S)$; hence $B\setminus S\subseteq N_{H[B]}(S)$ and
$S\cup N_{H[B]}(S)=B$.
\end{proof}

\begin{theorem}[Inductive reduction for ideal]\label{thm:ideal-reduction}
Let $F,G$ be AFs over a common finite universe with
$F\equiv^{\mathrm{nd}}_{\mathrm{ideal}}G$, and suppose the residual statement of
Conjecture~\ref{conj:hard} holds for $F,G$, namely:
\begin{quote}
$(\ast)$\quad for every $B$ and every maximal conflict-free set $S$ of $F[B]$,
$\;S\in\mathrm{adm}(F[B])\iff S\in\mathrm{adm}(G[B])$.
\end{quote}
Then $F\equiv^{\mathrm{nd}}_{\mathrm{adm}}G$.  Moreover, $(\ast)$ is the
\emph{only} obstruction: the equivalence $S\in\mathrm{adm}(F[B])\iff
S\in\mathrm{adm}(G[B])$ holds unconditionally for every $S$ that is not maximal
conflict-free, given $(\ast)$ on the strictly smaller frameworks below $B$.
\end{theorem}

\begin{proof}
Since $\equiv^{\mathrm{nd}}_{\mathrm{ideal}}$ refines
$\equiv^{\mathrm{nd}}_{\mathrm{cf}}$ (Lemma~\ref{lem:cf-recover} and
Remark~\ref{rem:cf-others}), we have $\mathrm{cf}(F[B])=\mathrm{cf}(G[B])$ for all
$B$; hence $F[B]$ and $G[B]$ share the same conflict-free sets, the same
conflict neighbourhoods $N(\cdot)$, the same closures $S\cup N(S)$, and the same
maximal conflict-free sets.  We use this throughout.

We prove, by induction on $|B|$, the statement
\[
\mathcal Q(B):\qquad \mathrm{adm}(F[B])=\mathrm{adm}(G[B]).
\]

\emph{Base case $|B|\le 1$.}  Every conflict-free subset of a framework on at
most one argument is maximal conflict-free (the only conflict-free sets are
$\emptyset$ and, if the single argument is not a self-attacker, the whole set,
which is then maximal).  Non-conflict-free sets are admissible in neither $F[B]$
nor $G[B]$.  Hence $\mathcal Q(B)$ follows directly from $(\ast)$.

\emph{Inductive step.}  Let $|B|\ge 2$ and assume $\mathcal Q(B')$ for all
$B'\subsetneq B$.  Fix $S\subseteq B$.  If $S$ is not conflict-free it is
admissible in neither framework, so we may assume $S\in\mathrm{cf}(F[B])=
\mathrm{cf}(G[B])$.

\emph{Case 1: $S$ is maximal conflict-free in $F[B]$ (equivalently in $G[B]$).}
Then $S\in\mathrm{adm}(F[B])\iff S\in\mathrm{adm}(G[B])$ by $(\ast)$.

\emph{Case 2: $S$ is conflict-free but not maximal.}  By Lemma~\ref{lem:maxcf},
$D:=S\cup N(S)$ satisfies $D\subsetneq B$.  By
Lemma~\ref{lem:conf-closure} applied to $F$ and to $G$,
\[
S\in\mathrm{adm}(F[B])\iff S\in\mathrm{adm}(F[D]),\qquad
S\in\mathrm{adm}(G[B])\iff S\in\mathrm{adm}(G[D]).
\]
By the inductive hypothesis $\mathcal Q(D)$ (as $|D|<|B|$),
$\mathrm{adm}(F[D])=\mathrm{adm}(G[D])$, so $S\in\mathrm{adm}(F[D])\iff
S\in\mathrm{adm}(G[D])$.  Combining the three displayed equivalences gives
$S\in\mathrm{adm}(F[B])\iff S\in\mathrm{adm}(G[B])$.

In both cases the membership of $S$ agrees, so $\mathrm{adm}(F[B])=
\mathrm{adm}(G[B])$, completing the induction.  As $B$ was arbitrary,
$F\equiv^{\mathrm{nd}}_{\mathrm{adm}}G$.

Finally, the ``moreover'' assertion is exactly Case~2: it used $(\ast)$ only
through the inductive hypothesis on the strictly smaller framework $H[D]$, and is
otherwise unconditional.  Thus the ideal--to--admissible problem reduces
\emph{exactly} to the maximal conflict-free statement $(\ast)$.
\end{proof}

\begin{remark}[Why the maximal--conflict-free case is genuinely directional]
\label{rem:directional}
For a maximal conflict-free $S$ of $H[B]$ one has $S\cup N(S)=B$
(Lemma~\ref{lem:maxcf}), so the conflict closure carries no proper-subset
structure to recurse on, and $S$ is admissible iff $S\in\mathrm{preferred}(H[B])$
(a maximal conflict-free admissible set is a maximal admissible set).  Whether
$S$ is admissible depends on the \emph{direction} of the conflicts between $S$
and $B\setminus S$ --- precisely the information that the conflict-free family
discards and that the ideal extension of $H[B]$ may hide when several preferred
extensions are present (e.g.\ a mutual attack $\{(s,b),(b,s)\}$ between $s\in S$
and $b\notin S$ makes $S$ admissible while leaving $\{b\}$-type rivals that can
collapse the ideal of $H[B]$).  This is why no single-restriction recovery of
admissibility from the ideal extension exists, and the residual case cannot be
reduced further by the present method.
\end{remark}

\begin{conjecture}[Residual maximal--conflict-free recovery]\label{conj:hard}
If $F\equiv^{\mathrm{nd}}_{\mathrm{ideal}}G$ then for every $B$ and every maximal
conflict-free set $S$ of $F[B]$ one has $S\in\mathrm{adm}(F[B])\iff
S\in\mathrm{adm}(G[B])$.  Equivalently, the maximal conflict-free admissible
(\,=\,preferred\,) sets of each restriction are determined by the family of ideal
extensions of all restrictions together with the conflict-free family.  We have
verified Conjecture~\ref{conj:hard} exhaustively over all $512$ AFs on the
$3$-element universe (where, by Theorem~\ref{thm:ideal-reduction}, it is
equivalent to $\equiv^{\mathrm{nd}}_{\mathrm{ideal}}=
\equiv^{\mathrm{nd}}_{\mathrm{adm}}$, both having $131$ classes) and by random
search on the $4$-element universe; a determinism check (the maximal-cf preferred
sets at $B$ are a function of the conflict-free family at $B$, the ideal
extension at $B$, the non-maximal-cf preferred sets at $B$, and the admissible
families of all proper restrictions of $B$) returned no violations.
\end{conjecture}

\begin{remark}\label{rem:cf-others}
Lemma~\ref{lem:cf-recover} also gives
$\equiv^{\mathrm{nd}}_{\mathrm{ideal}},
\equiv^{\mathrm{nd}}_{\mathrm{semi\text{-}stable}},
\equiv^{\mathrm{nd}}_{\mathrm{eager}}\subseteq
\equiv^{\mathrm{nd}}_{\mathrm{cf}}$: for ideal, $\mathrm{ideal}(H[S])=S$ iff
$S\in\mathrm{cf}(H[S])$ (when $S$ is conflict-free, $H[S]$ has no attacks, its
unique preferred extension is $S$ and hence its ideal extension is $S$; when $S$
is not conflict-free, $\mathrm{ideal}(H[S])\subsetneq S$).  An identical one-line
computation works for semi-stable and eager, since on the attack-free framework
$H[S]$ all these semantics return $\{S\}$.
\end{remark}

\section{Semi-stable, eager and stage}\label{sec:ss}

\begin{theorem}[Forward collapse for eager]\label{thm:eager-forward}
Define semi-stable deletion equivalence
$\equiv^{\mathrm{nd}}_{\mathrm{semi\text{-}stable}}$ as agreement of
$\mathrm{semi\text{-}stable}(\cdot[B])$ on all $B$.  Then
\[
\equiv^{\mathrm{nd}}_{\mathrm{semi\text{-}stable}}\cap
\equiv^{\mathrm{nd}}_{\mathrm{adm}}\ \subseteq\
\equiv^{\mathrm{nd}}_{\mathrm{eager}} .
\]
In particular, since (Theorem~\ref{thm:ss-adm}) semi-stable deletion
equivalence already implies admissible deletion equivalence, semi-stable
deletion equivalence implies eager deletion equivalence.
\end{theorem}

\begin{proof}
Fix $B$ and assume $\mathrm{semi\text{-}stable}(F[B])=
\mathrm{semi\text{-}stable}(G[B])$ and $\mathrm{adm}(F[B])=\mathrm{adm}(G[B])$.
The eager extension of $H[B]$ is, by definition, the $\subseteq$-largest
admissible set contained in the intersection $J_{H[B]}:=\bigcap
\mathrm{semi\text{-}stable}(H[B])$.  By Dung's fundamental lemma (the union of
all admissible subsets of a fixed set is again admissible), this largest set is
the union of all admissible sets contained in $J_{H[B]}$, hence a function of the
admissible family of $H[B]$ together with $J_{H[B]}$.  Both inputs coincide for
$F[B]$ and $G[B]$ (the second because the semi-stable extensions, hence their
intersection, coincide).  Thus $\mathrm{eager}(F[B])=\mathrm{eager}(G[B])$; as
$B$ was arbitrary, $F\equiv^{\mathrm{nd}}_{\mathrm{eager}}G$.
\end{proof}

\begin{theorem}[Semi-stable refines admissible]\label{thm:ss-adm}
$\equiv^{\mathrm{nd}}_{\mathrm{semi\text{-}stable}}\subseteq
\equiv^{\mathrm{nd}}_{\mathrm{adm}}$, and this inclusion is \emph{strict}: there
exist AFs that are admissible (equivalently preferred) deletion equivalent but
not semi-stable deletion equivalent.
\end{theorem}

\noindent\emph{Status.}  The inclusion and its strictness have been verified by
exhaustive enumeration on the $3$-element universe, where
$\equiv^{\mathrm{nd}}_{\mathrm{semi\text{-}stable}}$ has $137$ classes that
properly refine the $131$ classes of
$\equiv^{\mathrm{nd}}_{\mathrm{preferred}}=\equiv^{\mathrm{nd}}_{\mathrm{adm}}$;
exactly $3$ of the $131$ admissible classes split under semi-stable equivalence.
A self-contained proof of the inclusion would proceed as in
Lemma~\ref{lem:cf-recover}/\ref{lem:recover} but additionally recover the
\emph{range} datum $S\cup R(S)$ on every restriction from semi-stable
extensions; the range is precisely the extra information distinguishing the two
relations, and a closed-form recovery of it is the remaining technical step
(Section~\ref{sec:status}).  By Theorem~\ref{thm:eager-forward} the inclusion
$\equiv^{\mathrm{nd}}_{\mathrm{semi\text{-}stable}}\subseteq
\equiv^{\mathrm{nd}}_{\mathrm{eager}}$ holds once the present inclusion does; the
reverse inclusion (eager $\subseteq$ semi-stable) was likewise verified
exhaustively on the $3$-element universe, where both relations have $137$
classes.

\begin{proposition}[Four-way partition]\label{prop:partition}
On a universe of $3$ arguments the six semantics induce exactly four distinct
normal deletion equivalence relations, ordered by refinement as
\[
\underbrace{\{\mathrm{semi\text{-}stable},\mathrm{eager}\}}_{137}
\ \subsetneq\
\underbrace{\{\mathrm{preferred},\mathrm{ideal}\}}_{131},
\qquad
\{\mathrm{stage}\}_{92},
\qquad
\{\mathrm{naive}\}_{18},
\]
the subscripts being the number of classes; stage and naive are incomparable
with the admissibility chain in general.
\end{proposition}

\begin{proof}
The coincidences $\{\mathrm{preferred}=\mathrm{ideal}=\mathrm{adm}\}$ and
$\{\mathrm{semi\text{-}stable}=\mathrm{eager}\}$ and the strict refinement
between the two groups are Corollary~\ref{cor:pref-ideal},
Theorems~\ref{thm:eager-forward}--\ref{thm:ss-adm}; the closed form for naive is
Theorem~\ref{thm:naive}.  The exact class counts ($137,131,92,18$) and the
incomparabilities are obtained by exhaustive computation over all $512$ AFs on
the $3$-element universe.  Distinctness of the four groups is witnessed, e.g., by
\[
R_1=\{(a,c),(c,a),(b,b)\},\qquad
R_2=\{(a,c),(c,a),(a,b),(b,b)\},
\]
which are preferred- (hence admissible- and ideal-) equivalent but not
semi-stable equivalent; self-attacker examples separate stage from preferred and
from semi-stable.
\end{proof}

\section{Summary and open points}\label{sec:status}

\paragraph{Established rigorously (self-contained proofs).}
\begin{enumerate}
\item (Lemma~\ref{lem:reform}) For every semantics, normal deletion equivalence
is exactly equality of $\sigma$ on \emph{all} induced subframeworks; in
particular $A=A'$ is forced.
\item (Theorem~\ref{thm:naive}) \textbf{Naive.} $F\equiv^{\mathrm{nd}}_
{\mathrm{naive}}G$ iff $F,G$ have the same self-attackers and the same symmetric
conflict relation among non-self-attackers.  The naive deletion kernel is the
symmetric framework $(A,R^{\mathrm n})$.  Complete syntactic characterization.
\item (Lemma~\ref{lem:cf-recover}) \textbf{Conflict-free recovery.}
$S\in\mathrm{cf}(H[B])\iff S\in\mathrm{preferred}(H[S])$; hence
$\equiv^{\mathrm{nd}}_{\mathrm{preferred}}$ (and ideal, semi-stable, eager)
refine $\equiv^{\mathrm{nd}}_{\mathrm{cf}}$.
\item (Lemma~\ref{lem:recover}) \textbf{Admissibility recovery.}
$S\in\mathrm{adm}(H[B])\iff S\in\mathrm{preferred}(H[\,S\cup N_{H[B]}(S)\,])$,
where the conflict neighbourhood $N_{H[B]}(S)$ depends only on the conflict-free
family.
\item (Theorem~\ref{thm:pref-adm}) \textbf{Preferred $=$ admissible.}
$\equiv^{\mathrm{nd}}_{\mathrm{preferred}}=\equiv^{\mathrm{nd}}_{\mathrm{adm}}$,
proved \emph{both} directions from the two recovery lemmas.
\item (Theorem~\ref{thm:adm-forward}, Theorem~\ref{thm:eager-forward})
\textbf{Forward collapses.} Admissible deletion equivalence implies preferred and
ideal deletion equivalence; semi-stable $\cap$ admissible deletion equivalence
implies eager deletion equivalence.
\end{enumerate}

\paragraph{Established computationally (rigorous proof not enclosed).}
\begin{enumerate}
\item[(7)] (Corollary~\ref{cor:pref-ideal}) $\equiv^{\mathrm{nd}}_{\mathrm{ideal}}
\subseteq\equiv^{\mathrm{nd}}_{\mathrm{adm}}$ (so ideal $=$ preferred), verified
exhaustively on $3$ arguments and by random search on $4$.
\item[(8)] (Theorem~\ref{thm:ss-adm}) $\equiv^{\mathrm{nd}}_{\mathrm{semi\text{-}
stable}}\subseteq\equiv^{\mathrm{nd}}_{\mathrm{adm}}$ and
$\equiv^{\mathrm{nd}}_{\mathrm{eager}}\subseteq
\equiv^{\mathrm{nd}}_{\mathrm{semi\text{-}stable}}$, verified exhaustively on $3$
arguments (giving semi-stable $=$ eager, strictly finer than preferred) and by
random search on $4$.
\item[(9)] (Proposition~\ref{prop:partition}) The exact class counts
$137/131/92/18$ and the four-way partition with refinement
$\{\mathrm{semi\text{-}stable},\mathrm{eager}\}\subsetneq
\{\mathrm{preferred},\mathrm{ideal}\}$.
\end{enumerate}

\paragraph{Remaining genuine gaps.}
\begin{itemize}
\item \textbf{Ideal $\Rightarrow$ admissible (now reduced to one sharp case).}
We have \emph{proved} (Theorem~\ref{thm:ideal-reduction}, via
Lemmas~\ref{lem:conf-closure} and~\ref{lem:maxcf}) that ideal deletion
equivalence implies admissible deletion equivalence \emph{for every set that is
not maximal conflict-free}, by an induction on $|B|$ that reduces admissibility
of $S$ in $H[B]$ to admissibility of $S$ in the proper restriction
$H[S\cup N_{H[B]}(S)]$.  The implication $\equiv^{\mathrm{nd}}_{\mathrm{ideal}}
\subseteq\equiv^{\mathrm{nd}}_{\mathrm{adm}}$ is therefore equivalent to the
single residual statement (Conjecture~\ref{conj:hard}) that the
\emph{maximal conflict-free} admissible (= preferred) sets of each restriction
are recoverable from the ideal extensions of all restrictions and the
conflict-free family.  Remark~\ref{rem:directional} explains why this residual
case is genuinely directional and admits no single-restriction recovery, so it
cannot be reduced further by the present method; it has been verified
exhaustively on $3$ arguments and by random search on $4$.

\item \textbf{Semi-stable / eager $\Rightarrow$ admissible.}  Same recovery
strategy, plus recovery of the \emph{range} $S\cup R(S)$ of each set on each
restriction; the range is exactly the extra datum that makes semi-stable
deletion equivalence strictly finer than preferred deletion equivalence.
\item \textbf{Closed-form directed kernels.}  For the three directed relations
$\{$preferred, ideal$\}$, $\{$semi-stable, eager$\}$ and $\{$stage$\}$ a
$\subseteq$-minimal canonical attack relation exists on each finite class
(verified computationally), but it is \emph{not} given by any purely local rule
on the four edges $\{(a,b),(b,a),(a,a),(b,b)\}$ of a pair: whether an incoming
attack to a self-attacker is redundant is a non-local conflict-reachability
condition.  The classical source-side admissible-kernel reduction is sound but
not complete for these relations.
\end{itemize}

\paragraph{Conclusion.}  Naive deletion equivalence is completely and
syntactically characterized (Theorem~\ref{thm:naive}).  For the
admissibility-based semantics we have proved that \emph{preferred deletion
equivalence coincides exactly with admissible deletion equivalence}
(Theorem~\ref{thm:pref-adm}) via two framework-independent recovery lemmas.  For
the ideal semantics we have proved (Theorem~\ref{thm:ideal-reduction}) that ideal
deletion equivalence implies admissible deletion equivalence \emph{for every
non--maximal-conflict-free set}, reducing the equality
$\equiv^{\mathrm{nd}}_{\mathrm{ideal}}=\equiv^{\mathrm{nd}}_{\mathrm{adm}}$ to a
single sharply isolated residual case (Conjecture~\ref{conj:hard}: recovery of
the maximal conflict-free admissible sets), which Remark~\ref{rem:directional}
shows is genuinely directional and not amenable to single-restriction recovery.
Semi-stable and eager form a strictly finer relation obtained by additionally
tracking ranges.  The residual maximal-cf recovery for ideal, the reverse
inclusions for semi-stable/eager, and the closed-form directed kernels remain the
outstanding technical problems, confirmed computationally on small universes.

\end{document}
